\documentclass[11pt]{article}
\usepackage[margin=1in]{geometry}
\usepackage{amsmath,amssymb,mathtools}
\usepackage{fontspec,unicode-math}
\setmainfont{Latin Modern Roman}
\setsansfont{Latin Modern Sans}
\setmathfont{Latin Modern Math}
\usepackage{xurl,hyperref}
\hypersetup{colorlinks=true,urlcolor=blue,linkcolor=blue}
\newcommand{\sref}[2]{\hyperlink{source:#1:#2}{[S#2]}}
\newcommand{\cataloguescope}[1]{\paragraph{Recorded mathematical question.}#1}
\begin{document}
% UnsolvedMath ID 10000056; source catalogue code AMR-099-0056
\setcounter{section}{10000055}
\section{Recurrence under square-root separation limits}\label{10000056}
{\small\sffamily UnsolvedMath ID 10000056 · Graph Theory}
\subsection{Definitions and mathematical statement}

\paragraph{Shared notation.}$\mathbb N=\{1,2,\ldots\}$, and $\mathbb Z,\mathbb Q,\mathbb R,\mathbb C$ denote the integers, rationals, reals and complex numbers. Any other conventions are specified locally.


Let $G_k$ be finite connected graphs with a common degree bound. Assume that for every nonempty vertex set $S\subset V(G_k)$ one can delete at most $|S|^{1/2}$ vertices so that every connected component of the induced graph on the remaining vertices of $S$ has at most $|S|/2$ vertices. Root $G_k$ at a uniform random vertex $o_k$. Local convergence to a random rooted graph $(G,o)$ means convergence in law of every fixed-radius induced rooted neighborhood. A graph is recurrent if simple random walk returns to its starting vertex infinitely often almost surely.
\paragraph{Statement recorded in the supplied source.}
If $(G_k,o_k)$ converges locally to $(G,o)$, is $G$ recurrent almost surely? \sref{10000056}{2}
\subsection{Short English statement}
Does square-root separation, together with bounded degrees, force recurrence in a random local limit?
\subsection{Sources}
\begin{description}
\item[\hypertarget{source:10000056:1}{S1}] ULAM.AI and the UnsolvedMath Contributors, \emph{UnsolvedMath: A Curated Collection of Open Mathematics Problems}, record 10000056 (AMR-099-0056). CC BY 4.0. \url{https://huggingface.co/datasets/ulamai/UnsolvedMath}.
\item[\hypertarget{source:10000056:2}{S2}] Itai Benjamini, \emph{Euclidean vs. Graph Metric}, version dated 16 July 2012; original author manuscript preserved in Arquivo.pt capture of 31 December 2020, Section 2, definition of square-root separation and Question 2.2, pp. 2--3. \url{https://arquivo.pt/noFrame/replay/20201231041538id_/http://www.wisdom.weizmann.ac.il/~itai/erd100.pdf}.
\end{description}
\label{end:10000056}
\end{document}
